\documentclass{beamer}
\usepackage[utf8]{inputenc}
\usepackage{quantikz}
\usepackage{amsmath}
\usepackage{braket}

\title{Quantum Phase Estimation (QPE)}
\author{Morten Hjorth-Jensen}
\date{Slides from 2007}

\begin{document}

% Title Slide
\begin{frame}
    \titlepage
\end{frame}

% Outline Slide
\begin{frame}{Outline}
    \tableofcontents
\end{frame}

% Introduction Slide
\section{Introduction}
\begin{frame}{What is Quantum Phase Estimation?}
    \begin{itemize}
        \item Quantum Phase Estimation (QPE) is a fundamental quantum algorithm to estimate the eigenphase $\phi$ of a unitary operator $U$.
        \item Given an eigenstate $\ket{\psi}$ such that $U\ket{\psi} = e^{2\pi i \phi}\ket{\psi}$, QPE estimates $\phi$ with high probability.
        \item It is a crucial subroutine in algorithms like Shor's and quantum simulations.
    \end{itemize}
\end{frame}

% Mathematical Background
\section{Mathematical Background}
\begin{frame}{Mathematical Background}
    \begin{itemize}
        \item Consider a unitary operator $U$ with eigenstate $\ket{\psi}$:
        \[ U\ket{\psi} = e^{2\pi i \phi} \ket{\psi}, \quad \phi \in [0,1). \]
        \item Goal: Estimate $\phi$ to $n$ bits of precision.
        \item Quantum Fourier Transform (QFT) plays a key role in extracting phase information.
    \end{itemize}
\end{frame}

% Circuit Structure
\section{Circuit Structure}
\begin{frame}{Quantum Circuit for QPE}
    \begin{figure}
        \centering
        \begin{quantikz}
            \lstick{$\ket{0}^{\otimes n}$} & \qwbundle{n} & \gate[2]{QFT^\dagger} & \meter{} \\
            \lstick{$\ket{\psi}$} & \gate{U^{2^0}} & \gate{U^{2^1}} & \qw & \qw
        \end{quantikz}
        \caption{Quantum Phase Estimation Circuit with $n$ qubits.}
    \end{figure}
\end{frame}

% Derivation of the Algorithm
\section{Derivation of the Algorithm}
\begin{frame}{Step 1: Initialization}
    \begin{itemize}
        \item The system starts in $\ket{0}^{\otimes n} \otimes \ket{\psi}$.
        \item Apply Hadamard gates to the first $n$ qubits:
        \[ \frac{1}{2^{n/2}} \sum_{k=0}^{2^n -1} \ket{k} \otimes \ket{\psi}. \]
    \end{itemize}
\end{frame}

\begin{frame}{Step 2: Controlled Unitary Operations}
    \begin{itemize}
        \item Apply controlled-$U^{2^j}$ operations to create phase kickbacks:
        \[ \frac{1}{2^{n/2}} \sum_{k=0}^{2^n -1} \ket{k} \otimes U^k\ket{\psi}. \]
        \item For eigenstate $\ket{\psi}$ with eigenvalue $e^{2\pi i \phi}$:
        \[ U^k\ket{\psi} = e^{2\pi i \phi k} \ket{\psi}. \]
    \end{itemize}
\end{frame}

\begin{frame}{Step 3: Quantum Fourier Transform (QFT)}
    \begin{itemize}
        \item QFT on the first $n$ qubits transforms the state:
        \[ \frac{1}{2^{n/2}} \sum_{k=0}^{2^n -1} e^{2\pi i \phi k} \ket{k}. \]
        \item After QFT, we measure the most probable state representing $\phi$.
    \end{itemize}
\end{frame}

% Applications
\section{Applications}
\begin{frame}{Applications of QPE}
    \begin{itemize}
        \item Estimating eigenvalues of Hermitian operators.
        \item Integral part of Shor's algorithm for integer factorization.
        \item Quantum simulation for molecular energy estimation.
    \end{itemize}
\end{frame}

% Conclusion
\section{Conclusion}
\begin{frame}{Conclusion}
    \begin{itemize}
        \item QPE is a powerful quantum algorithm for phase estimation.
        \item Leverages QFT and controlled unitary operations.
        \item Forms a foundational block for many quantum algorithms.
    \end{itemize}
\end{frame}

% References
\begin{frame}{References}
    \begin{itemize}
        \item M. A. Nielsen and I. L. Chuang, \textit{Quantum Computation and Quantum Information}.
        \item Quantum algorithms lecture notes.
    \end{itemize}
\end{frame}

\end{document}
